\begin{table}[!tbp]\centering
\caption{Selection bounds and statistical power}\label{tab:leemde}
\begin{threeparttable}
\footnotesize
\begin{tabular}{lcc}
\toprule
\multicolumn{3}{l}{\emph{Panel A: Lee-type trimming bounds, unconditional 2$\times$2 log-wage DiD}}\\
\midrule
Selection DiD (wage-sample inclusion) & -0.021 & \\
Trimming fraction & 0.076 & \\
Untrimmed DiD & -0.005 & \\
Bounds & [-0.099, 0.101] & \\
\midrule
\multicolumn{3}{l}{\emph{Panel B: minimum detectable effects (80\% power, 5\% size)}}\\
Outcome & Estimate & MDE \\ \midrule
Log real hourly wage & -0.021 & 0.034 \\
Log real monthly earnings & -0.010 & 0.038 \\
Employed (working age) & -0.025 & 0.016 \\
Labour force participation & -0.008 & 0.012 \\
Formal employment & -0.046 & 0.021 \\
Weekly hours (main job) & 0.556 & 0.698 \\
Full-time (>=35h) & -0.013 & 0.013 \\
Self-employed & 0.036 & 0.017 \\
Paid below the MW & -0.006 & 0.021 \\
\bottomrule
\end{tabular}
\begin{tablenotes}\footnotesize
\item Notes: Panel A: bounds on the unconditional 2$\times$2 log hourly wage DiD following the trimming logic of \citet{lee_2009}: the reform reduced the probability that a low-education worker appears in the wage-earner sample (selection DiD), so the low$\times$post cell is trimmed from above and below by the implied fraction of its baseline rate. Panel B: minimum detectable effect $=2.8\times$ the department-clustered standard error of the baseline DiD.
\end{tablenotes}
\end{threeparttable}\end{table}
